\cotefile{79}
\scanpage{113}\mspage{57}

générateurs satisfaisant

\[
\rho_0,\rho_1,\rho_\infty\quad\mid\quad \rho_0^n=\rho_1^2=\rho_\infty^3=1
\]

% slanted
\margin{\uline{NB} $n$ est exactement l'ordre de $\rho_0$ dans $G_0^+$.}

(rotations autour de $s_0$, du centre de $\{s_0,t_0\}$, et du centre de la face $\{s_0,t_0,u_0\}$), et que $G_{0\,s_0}^+$ est le sous-groupe engendré par $\rho_0$, la condition $(T_0\,3')$ est \deleted{donc} une traduction de $T_0\,3$ en termes de théorie des groupes pure. \deleted{\unsure{Dans}}

Dans les cas qui nous intéressent, où $P_0$ est régulier = en \deleted{\unsure{tant}} que polyèdre non orienté, on aura pour \deleted{\unsure{$z$|$Z$} \unsure{de}} $Z'$ commutant = à $G_0$ tout entier, non seulement à $G_0^+$, et on pourra, dans l'interprétation précédente, choisir un $g\in G_0$ -- p.ex. on pourra prendre la \deleted{symétrie}\inserted{réflexion} $\tau_0$ \deleted{\unsure{pour}} qui échange $s_0$ et $t_0$, la formule $(3\deleted{\unsure{3}}4)$ devient

\[
g\cdot\operatorname{int}(\tau_0)(gz)\ \in\ \langle\rho_0\rangle \tag*{$(34')$}
\]

\margin{sous-groupe engendré par $\rho_0$}

pour $z\in\langle\rho_0\rangle$ convenable (dépendant du choix de $g\in\operatorname{Norm}_{\deleted{\unsure{G_0^+}}}(\langle\rho_0\rangle)\unsure{)}\bigr]\unsure{]}$
% archive sheet number: 112
\scanpage{114}

Considérons maintenant, pour $\lambda\in Z'$

\[
\begin{aligned}
 A_\lambda&=\left\{a\in\unsure{\mathbb P_2|\mathcal P_2|\mathfrak P_2}(S)\ \middle|\ 
 \begin{gathered}\exists(s,t)\in S^2,\ \text{avec}\\
 \{s,t\}\in A_0,\quad a=\{s,\lambda.t\}\end{gathered}\right\}\\
 &=\bigl\{\{s,\lambda t\}\mid
 \underbrace{(s,t)\in S^2,\ \text{tq. }\{s,t\}\in A_0}_{\substack{\text{\uline{NB} cette condition implique}\\
 \text{bien par (31) que }s\ne\lambda t\\
 \text{donc }\{s,\lambda t\}\in\mathbb P_2(S)}}\bigr\}
\end{aligned}\tag*{$(35)$}
\]

On a donc \inserted{(pour $\lambda$ choisi!)} une \deleted{\unsure{$A$|$A_\lambda$}} partie de $\mathbb P_2(S)$, définissant sur $S$ une structure de graphe strict, notée $\Gamma_\lambda$, qui est liée à la structure de graphe $\Gamma_0=(S,A_0)$ associée à \deleted{\unsure{$P_0$}} $P_0=(S,A_0,F_0,\ldots)$, via $\lambda$, par le fait que $\lambda$ induit une bijection

\[
\underset{\substack{\text{étoile de}\\s\text{ pour }\Gamma_0}}{A_0(s)}
 \xrightarrow{\ \sim\ }
\underset{\substack{\text{étoile de}\\s\text{ pour }\Gamma_\lambda}}{A_\lambda(s)}
\tag*{$(36)$}
\]

puisqu'on aura, pour $\unsure{\{s,t\}|(s,t)}\in S^2$, \unsure{)}

\[
\underset{\text{i.e. }t\in A_0(s)}{\{s,t\}\in A_0}\quad\Longleftrightarrow\quad
\underset{\lambda t\in A_\lambda(s)}{\{s,\lambda t\}\in A_\lambda\ \text{i.e}}\ ,
\tag*{$(37)$}
\]

ce qui est essentiellement la définition (35) de $A_\lambda$, à cela près que celle ci signifie que $\underset{\substack{=\{\lambda t,s\}\\=\lambda t,\lambda(\lambda^{-1}s)}}{\{s,\lambda t\}}\in A_\lambda$ ssi \deleted{\unsure{ou}} \unsure{--} $\{s,t\}\in A_\lambda$ \uline{ou} $\{\lambda^{-1}s,\lambda t\}\in A_\lambda$, mais par $T_0\,2$ les deux conditions sont équivalentes
% archive sheet number: 113
\scanpage{115}\mspage{58}

On voit donc que pour $\forall\,s\in S$, l'ordre du \deleted{\unsure{$s$}} sommet $s$ pour le graphe $A_\lambda$ est \deleted{\unsure{aussi}} égal à son ordre \inserted{$n$} pour $A_0$. Notons aussi que par transport des structures, $G_0^+$ \deleted{opère} transforme $A_\lambda$ en lui-=, i.e. $G_0^+$ opère par autom. du graphe $\Gamma_\lambda$. \deleted{ainsi} \deleted{qui est donc} On voit aussitôt qu'il est \inserted{simplement} transitif sur l'ens. des arêtes orientées de \deleted{\unsure{$G$}}$\Gamma_\lambda$, puisqu'il en est ainsi pour $\Gamma_0$.

On veut exprimer que $\Gamma_\lambda$ définit une structure de polyèdre régulier triangulé isomorphe à $P_0$. On va supposer pour ceci, d'abord, que \deleted{\unsure{le}}la structure de graphe $\Gamma_0\inserted{=(S_0,A_0)}$ détermine la structure de polyèdre triangulé $P_0=(S_0,F_0)$; \deleted{\unsure{p.ex.|i.e.}}

\[
(T_0\,4)\quad \deleted{\unsure{\text{on}}}\ \text{a}\quad F_0=\left\{f\in\unsure{\mathbb P_3|\mathcal P_3|\mathfrak P_3}(S)\ \middle|\ 
 \begin{gathered}\forall a\in\mathbb P_2(f),\ \text{on a}\\a\in A_0\end{gathered}\right\}
\tag*{$(38)$}
\]

% slanted
\margin{devrait venir avant $T_0\,1,T_0\,2,T_0\,3$ $\longleftarrow$}

(la structure de graphe détermine la structure triangulée)

J'ignore, s'il existe une structure de 2-polyèdre régulier \deleted{orientable} triangulé orientable, qui ne
% archive sheet number: 114
\scanpage{116}

soit ainsi décrit par sa structure de graphe déjà -- d'ailleurs même sans supposer la structure orientable ni régulière, ce n'est pas évident qu'il y ait un contre-exemple. \deleted{Mais si --} Mais si -- dans l'\uline{icosaèdre gauche}, (régulier, mais pas orientable), on  $S$ de cardinal 6, \deleted{\unsure{tte}} \deleted{\unsure{tout|toute}} $a\in\unsure{\mathbb P_2|\mathcal P_2|\mathfrak P_2}(S)$ est une arête (il y en a 15), mais seulement la moitié des 20 éléments $f\in\mathbb P_3(S)$ sont des faces\ldots.

% slanted
\margin{\uline{Digression sur relation entre $A_0$ et $F_0$}}

On peut essayer de voir, pour une structure de graphe \inserted{strict}, $\Gamma_0=(S_{\unsure{0|2}},A_0)$, à quelles conditions l'ens $F_0\subset\mathbb P_3(S_{\unsure{0}})$ qu'il \deleted{\unsure{en}} définit est une structure de 2-polyèdre. On trouve \deleted{\unsure{ce|ceci}}

\[
\begin{aligned}
&\deleted{\unsure{(40)}}\\
&A1)\quad \forall\,\underset{\substack{\shortparallel\\\{s,t\}}}{a}\in A_0,\ 
 \text{il existe exactement deux éléments}\\
&\qquad u\in S,\ \text{tels que l'on ait}\\
&\qquad \{s,u\},\{t,u\}\in A_0\quad \text{i.e. adjacents}\\
&\qquad \text{à la fois aux deux extrémités}\\
&\qquad \text{de l'arête }a\quad\text{i.e. }A_0(s)\cap A_0(t)\text{ est}\\
&\qquad \text{de cardinal }2.
\end{aligned}\tag*{$(39)$}
\]

\margin{\uline{NB} le fait qu'il existe au, moins \uline{un} tel $u$ indique déjà que $A_0$ se déduit de $F_0$ par $A_0=\{a\in\mathbb P_2(S)\mid\exists f\in F_{\unsure{3|0}}\ f\supset a\}$}

\[
\begin{aligned}
&A2)\quad \forall s\in S,\quad\text{la structure de contour dont l'ens.}\\
&\qquad\text{des sommets est }\inserted{A_0(s)=}\operatorname{Et}_{\Gamma_0}(s)\quad\text{(ens.}\\
&\qquad\text{des sommets adjacents), et où}\\
&\qquad\text{les arêtes sont les } a\in\mathbb P_2(A_0(s))\\
&\qquad\text{telles que }a\in A_0,\quad\text{est \uline{connexe}}\\
&\qquad\text{non vide \deleted{(\uline{NB} si elle est non v.}}
\end{aligned}
\]

\margin{\uline{NB} Si $A_0(s)\ne\varnothing$, il résulte déjà de $A1$ que $\operatorname{Card}A_0(s)\geqslant3$}
% archive sheet number: 115
\scanpage{117}\mspage{59}

Ceci vu, on peut se poser la question si $A_\lambda$, déduit de $A_0$ via $\lambda$ comme dit dans (35), définit bien une structure de graphe satisfaisant à $A1)$ $A2)$, i.e. permettant de définir une structure 2-polyédrale (qui sera automatiquement régulière \deleted{\unsure{)}}, on trouve en $G_0^+$ \deleted{\unsure{sur|formant}} un groupe d'automorphismes transitif sur l'ens des arêtes orientées du dit polyèdre\ldots).

Considérons d'abord $A1$, soit donc \deleted{$a\in\unsure{\mathbb P_2|\mathcal P_2|\mathfrak P_2}\unsure{(|(S)}$} $a\in A_\lambda$, on peut écrire

\[
a=\{s,\lambda t\}\qquad\text{avec }\{s,t\}\in A_0
\]

Soit $v\in S$, donc $v$ de la forme $\lambda u$, on aura donc

\[
\begin{aligned}
\deleted{\{s,\lambda t,v\}}&\\[-2pt]
{\underset{\substack{\text{i.e. }\{s,v\}\in A_\lambda\\\{\lambda t,v\}\in A_\lambda}}{a\cup\{v\}\in F_\lambda}}
 &\Longleftrightarrow\quad
 \begin{gathered}\{s,\deleted{\unsure{v|u}}\lambda u\}\in A_{\deleted{\unsure{0}}\lambda}\\
 \{\lambda t,\lambda u\}\in A_\lambda\end{gathered}\\
 &\Longleftrightarrow\quad
 \begin{gathered}\{s,u\}\in A_0\\\{t,\lambda u\}\in A_0\end{gathered}
\end{aligned}
\]

donc la condition $A1$ pour $\deleted{\unsure{A}}\Gamma_\lambda$ s'écrit

\[
\begin{aligned}
&\deleted{\text{Pour }(s,t)\in\unsure{S^2|A_0^2},\ \text{tel que }\{s,t\}\in}\\
&\text{Pour deux sommets }s,t\text{ adjacents pour }P_0,\ \inserted{\text{et }\lambda\in Z',}\text{ il existe}\\
&\text{exactement deux sommets }u\text{ tels que }\deleted{\text{l'on ait}}\\
&u\text{ soit adjacent à }s,\quad\lambda u\text{ adjacent à }t,\ \text{i.e.}\\
&\operatorname{Card}\ \lambda.A_0(s)\cap A_0(t)=2\qquad \text{si }
 \begin{cases}\{s,t\}\in A_0\\\lambda\in Z'\end{cases}
\end{aligned}\tag*{$(40)$}
\]
% archive sheet number: 116
\scanpage{118}

Cette condition n'a pas l'air tautologique !

D'autre part, la condition de connexité pour $A_\lambda(s)$ peut s'examiner à l'aide de la bijection canonique (36) induite par $\lambda$

\[
A_0(s)\xrightarrow{\ \sim\ } A_\lambda(s)
\]

on note que $\{\lambda t,\lambda t'\}\in A_\lambda$ \unsure{sss|ssi} $\{t,\lambda t'\}\in A_0$ -- -- pourquoi cette relation \inserted{symétrique} entre deux éléments $t,t'$ de $A_0(s)$ serait elle ``connexe'' ?

% slanted
\margin{on voit bien, sur l'exemple de la double pyramide (sur un $n$-polygone avec $n$ pair $\geqslant6$, $n\not\equiv0\ (4)$) que cette condition n'est \uline{pas} tautologique !}

En fait, on va \unsure{apaiser|applanir} ces perplexités par un axiome qui, contrairement à l'axiome ``par acquit de \unsure{conscient|conscience}'' $T_0\,4$, a l'air des plus substantiels

\[
\begin{aligned}
(T_0\,5)\qquad &\forall\lambda\in Z',\quad\text{la structure de graphe}\\
&\Gamma_\lambda=(S,A_\lambda)\quad(\text{cf }35)\quad\text{est \uline{isomorphe}}\\
&\text{à }\Gamma_0=(S,A_0),
\end{aligned}\tag*{$(41)$}
\]

en d'autres termes, on demande qu'il existe une permutation $g=g_\lambda\in\unsure{\mathfrak S_S|G_S}$ telle que l'on ait

\[
\forall s,t\in S,\ \text{on a}\quad \{s,t\}\in A_0\Longleftrightarrow
 \underbrace{\{gs,\lambda^{-1}gt\}\in A_0}_{\substack{\text{ce qui exprime}\\\{gs,gt\}\in A_\lambda}}
\tag*{$(42)$}
\]
% archive sheet number: 117
\scanpage{119}\mspage{60}

Du coup il est clair que $A_\lambda$ définit une structure de 2-polyèdre régulier et dont et \deleted{tout}

\[
P_\lambda=\bigl(S,\underset{\substack{\cap\\\unsure{\mathbb P_2|\mathcal P_2|\mathfrak P_2}(S)}}{A_\lambda},
 \underset{\substack{\cap\\\mathbb P_3(S)}}{F_\lambda},\ \text{relations d'incidence évidentes}\bigr)
\tag*{$(43)$}
\]

En fait, on a vu que $G_0^+$ opère par automorphismes de $\Gamma_\lambda$, donc il opère par automorphismes de $P_\lambda$ -- on devine que ce sont des automorphismes directs -- \margin{à y revenir} mais on voudrait que $G_0$ aussi opère \deleted{sur} par $P_\lambda$-automorphismes, pour ceci on va supposer

\[
\begin{aligned}
(T_0\,6)\quad &Z'\text{ commute (non seulement à}\\
&G_0^+,\text{ mais) à tout le groupe }G_0,
\end{aligned}\tag*{$(\deleted{\unsure{42}}\unsure{42|43|44})$}
\]

% slanted
\margin{on pourrait aussi continuer sans cet axiome, en travaillant avec le sous-groupe de $Z'$ commutant à $G_0$ tt entier -- en ignorant \unsure{=|--} l'orientabilité de $P_0$ i.e. laissant tomber l'axiome $T_0\,2$.}

qui implique que $G_0$ opère sur $S$ par $\Gamma_\lambda$-automorphismes, donc par $P_\lambda$-automorphismes. \deleted{On} J'aimerais vérifier que l'hom. canonique injectif

\[
\deleted{\unsure{(44)}}\qquad G_0\lhook\joinrel\longrightarrow\operatorname{Aut}(P_\lambda)\overset{\text{déf}}{=}G_\lambda
\]

est un isom -- comme $P_\lambda\simeq P_0$, on a $G_\lambda\simeq G_0$, et l'isomorphie est évidente si
% archive sheet number: 118
\scanpage{120}

$G_0$ est fini i.e. $S$ fini. Sinon ce n'est pas tellement clair sur les définitions. On peut quand = dire que \deleted{pour} $G_0^+$ est simplement transitif sur l'ens. $\vec A_\lambda$ des arêtes orientées de \deleted{$A$} $P_\lambda$, donc \deleted{\unsure{si on a|si on savait}} pour montrer que $G_0$ est transitif (donc \deleted{deux repères $\unsure{r,r'}$ de $P_\lambda$,} simpl. transitif) sur l'ens des repères de $A_{\unsure{\lambda|0}}$, \deleted{$r_0=(s_0,t_0,u_0)\qquad r_1=(s_1,t_1,u_1)$} \deleted{il} il suffit de prouver que pour \deleted{\unsure{tt}} un repère $(s_0,t_0,u_0)$\inserted{(donc pour \unsure{tt|$\forall$} repère)}, il existe un $u\in G_0$ \deleted{tel} qui fixe $s_0,t_0$, et \uline{non} $u_0$, \deleted{Mais s'il}\inserted{i.e.} qui n'est pas l'identité. Mais s'il n'en était ainsi, cela signifierait que $G_0$ est \uline{simplement} transitif sur l'ens $\vec A_\lambda$ (le stabilisateur d'un $\vec a\inserted{=(s_0,t_0)}\in\vec A_\lambda$ dans $G$ est réduit à l'élément neutre), ce qui contredit le fait que $G_0^+$ est simplement transitif et \deleted{\unsure{d'i. de}} $G_0^+\subsetneq G_0$. Donc on a bien un iso

\[
G_0\mathrel{\overset{\sim}{\rule{2cm}{0.4pt}}}\operatorname{Aut}(P_\lambda)=G_{\lambda\deleted{\unsure{\lambda|?}}},
\tag*{$(43)$}
\]

on aimerait prouver que cet iso. induit

\[
G_0^+\xrightarrow{\ \sim\ }G_\lambda^+\qquad ?
\tag*{$(44)$}
\]
% archive sheet number: 119
\scanpage{121}\mspage{61}

Cela se voit facilement si l'ordre des points de $P_0$ est impair, car alors $G_0^+$ se décrit intrinsèquement \deleted{comme} en termes de $G_{\unsure{0}}$, par le fait que c'est le \uline{seul} sous-groupe d'indice 2. Si par contre $n$ est pair, ce n'est pas chose évidente -- il y a \uline{trois} sous-groupes d'indice 2 possibles \inserted{a priori}, car $(\unsure{\mathbb G_{\infty\,3}|G_{\infty\,3}})_{\mathrm{ab}}\simeq\mathbb F_2\times\mathbb F_2$. Mais il n'y en a qu'un seul qui contienne $\rho_1^r$, donc la question est si $\rho_1^r$ (relatif à un repère $r\deleted{\unsure{\}}}=\{s_0,t_0,u_0\}$ de $P_0\}$ conserve l'orientation \deleted{dans} $P_\lambda$\ldots

% slanted
\margin{il faudra revenir là-dessus\ldots}

Je m'intéresse maintenant à l'ensemble des structures $P_\lambda$, \inserted{$\lambda\in Z'$.} \deleted{Notons d'abord} Les structures polyédrales triangulées qui satisfont aux axiomes $(T_0\,1)$ à $(T_0\,6)$ seront appelées (à défaut de mieux, et provisoirement) ``spéciales''. Je dirais que deux telles structures, \inserted{définies par des ens. d'arêtes $\subset\unsure{\mathbb P_2|\mathcal P_2|\mathfrak P_2}(S)$} $\deleted{\unsure{F_0}}A_0,\deleted{\unsure{F_1}}A_1$ sur un même ensemble de sommets \inserted{$S$} sont \uline{équivalentes}, il existe un $\lambda\in\unsure{\mathfrak S_S|G_S}$, commutant à $\operatorname{Aut}(S,F_0)$, de telle façon que l'on ait
% archive sheet number: 120
\scanpage{122}

\[
\forall(s,t)\in S,\ \text{on a}\qquad
 \{s,t\}\in A_0\Longleftrightarrow\{s,\lambda t\}\in A_1.
\tag*{$(45)$}
\]

On voit que cela implique que \deleted{\unsure{la}} les groupes d'autom. $G_0,G_1$ de $P_0,P_1$ sont égaux, et que par suite $\lambda$ \inserted{donc $\lambda^{-1}$} commute aussi à $G_{1}$, et que \deleted{\unsure{l'on|l'ensemble}} par suite $A_0$ est ``transformé'' de $A_1$ par l'élément $\lambda^{-1}$ du commutant de $G_1$ dans $\unsure{\mathfrak S_S|G_S}$ : la relation envisagée est donc symétrique, elle est évidemment réflexive, et on voit tout de suite qu'elle est également transitive : \deleted{elle} c'est donc bien une \uline{relation d'équivalence} entre structures polyédrales spéciales. Si \deleted{\unsure{$C$}} \inserted{$\Pi$} est une classe d'équivalence de telles structures, il lui correspond un groupe

% slanted
\margin{on note l'ens. $\Pi$ \unsure{indiciellement|radicalement}}

\[
\begin{aligned}
 G_0,\quad &\text{groupe d'automorphismes}\\
 &\text{commun aux structures}\\
 &\deleted{\unsure{F}}P_i\quad\deleted{\unsure{\in C}}\quad(i\in\Pi)
\end{aligned}\tag*{$(46)$}
\]

et un groupe

\[
 Z'=\operatorname{Comm}_{\mathfrak S_S}(G_0)\ ,
\tag*{$(47)$}
\]

enfin, on voit que le groupe $Z'$ opère sur $\Pi$, transitivement par définition
% archive sheet number: 121
\scanpage{123}\mspage{62}

de notre relation d'équivalence. $[\![$ C'est le moment peut-être, d'ailleurs, de noter la relation plausible

\[
Z'\text{ est un groupe \uline{commutatif}}
\tag*{$(48)$}
\]

% slanted
\margin{Je ne vais pas la poser en axiome pour \uline{l'instant} cf la démonstration bien \unsure{foireuse|foirée} suivante}

Car si $\lambda,\mu\in Z'$, et $(s,t)\in S^{\unsure{2}}$, on a

\diagram{\small\begin{tikzcd}[ampersand replacement=\&,column sep=small,row sep=large]
(s,t)\in A_0 \arrow[r,Leftrightarrow,"\lambda\in Z'"] \arrow[d,Leftrightarrow,"\lambda\mu\in Z'"'] \&
 \{\lambda s,\lambda^{-1}t\}\in A_0 \arrow[r,Leftrightarrow,"\mu\in Z'"] \&
 \{\mu\lambda s,\mu^{-1}\lambda^{-1}t\}\in A_0\\
 \{(\lambda\mu)s,(\lambda\mu)^{-1}t\}\in A_0\mathpunct{,} \& \& 
\end{tikzcd}}

donc en comparant les extrêmes, et posant $\mu^{-1}\lambda^{-1}t=(\lambda\mu)^{-1}t=u$

\[
\{u,(\lambda\mu)s\}\in A_0\Longleftrightarrow\{u,(\mu\lambda)s\}\in A_0
\]

$\forall s,u\in S$, ce qui s'écrit \deleted{\unsure{ce}}

\[
\deleted{\unsure{(49)}}\quad A_0((\lambda\mu)s)=A_0((\mu\lambda)(s))\ ,
\]

égalité des \uline{étoiles} de $(\lambda\mu)(s)$ et $(\mu\lambda)(s)$.

\margin{\deleted{\unsure{g}}}

Si donc on suppose que l'un sommet est déterminé par son étoile, on trouve bien (48) -- mais je ne vois pas \deleted{faire} reposer l'axiome explicite sur cette détermination un sommet par son étoile -- (qui est d'ailleurs \deleted{faux} pour l'octaèdre\ldots). $]\!]$

% slanted
\margin{c'est vrai du fait que $P_0$ est simpl. triangulé \unsure{connexe|convexe} dont \uline{tt} les sommets sont de \unsure{égal|même} ordre, à l'exception du seul octaèdre}

% slanted
\margin{cf bas de la page 22 pour un argument bizarre \unsure{en ce sens|réussi|aussi}\ldots}
% archive sheet number: 122
\scanpage{124}

La commutativité de $Z'$ finalement résulte déjà des axiomes $T_0\,3,T_0\,4$ ($T_0\,1$ et \uline{$T_0\,2$} sont \unsure{inutiles} !), pour un 2-polyèdre triangulé connexe, \deleted{\unsure{dans|dont}} \unsure{le cas} \uline{où} un sommet n'est pas déterminé par son étoile, alors le polyèdre est une double pyramide \inserted{(à base un $n$-gone ($n\geqslant3$))} \unsure{pour $n\ne4$ son} dont le groupe des automorphismes est \deleted{\unsure{le}}\inserted{pr.duit} du groupe diédral $\mathbb D_n$ \inserted{($\simeq G_0^+$)} par le groupe $\{1,\underline-\}$ engendré par l'antipodisme.

% slanted
\margin{sinon pour \unsure{déjà} introduire $Z'$ comme \unsure{commutant} de $G_0^+$ -- \unsure{alors qu'on aurait pu} le définir comme \unsure{celui} de $G_0$ tout \unsure{entier}\ldots}

Le commutant de $G_0^+$ \inserted{dans $\unsure{\mathfrak S_S|G_S}$} est réduit à : \deleted{$\unsure{\{1,\sigma\}|\{1,\varepsilon\}}\times\{1,\underline-\}$ pour $n$ impair (centre de} \inserted{$\mathfrak z\times\{1,\underline-\}$, où $\mathfrak z$ est le} centre de $G_0^+\simeq\mathbb D_n$ (réduit à $\{1\}$ si $n$ impair, et à $\mathbb Z/2\mathbb Z$ si $n$ pair) donc dans ce cas le commut\textsuperscript{ant}\ \unsure{=}\ est commutatif.

% slanted
\margin{le commutant de $G_0^+$ est \inserted{bien} égal à celui de $G_0$, et égal au centre de $G_0$\unsure{,}}

Reste le cas (plus intéressant pour nous) où $n=4$, i.e. celui où tous les sommets ont = ordre (c'est le cas de l'octaèdre. Dans \unsure{le|le cas} le commutant de $G_0^+$ \inserted{$\simeq\mathfrak S_4$} (égal du reste à celui de $G_0\simeq G_0^+$, et au centre de $G_0$) est réduit à $\{1,\underline-\}$. \uline{En tous cas $Z'$ est commutatif}.

% slanted
\margin{En tous cas, le commutant de $G_0^+$ \deleted{est égal au} centre de $G_0$, est, soit égal à $\{1,\underline-\}$ (cas $n$ impair ou $n=4$) soit à $D_2=\{1,\sigma\}\times\{1,\underline-\}$ (cas $n$ pair $\geqslant6$) -- \unsure{à ch.} un d'élément du commutant étant formé d'éléments de $G_0$ involutifs, satisfait trivialement $T_0\,3$. Mais dans le cas $n$ pair $\geqslant6$, la structure de graphe $A_\sigma$ déduite de $A_0$ en appliquant $\sigma$ (fixant les sommets $s,s'$ de la double pyramide, et transformant chaque sommet de $E$ en son opposé) ne correspond \unsure{,|qu'}à une structure \unsure{triangulée|de} \deleted{\unsure{de}} que si $n\equiv0\ (4)$}

Donc ds le cas symétrique [où tous les pts ont = ordre], on voit

\[
\left\{\begin{aligned}
a)&\ \text{ou bien tt \deleted{point}\inserted{sommet} est déterminé par son étoile,}\\
b)&\ \text{ou bien }P_0\text{ est un octaèdre, }\inserted{\text{(cas régulier)}}\
 Z'=\{1,\underline-\},\quad\underline-\\
 &\text{l'antipodisme, i.e. l'unique élément}\\
 &\text{central }\ne\mathrm{Id}\text{ de }G_0=\operatorname{Aut}(P_0)
 (\simeq\mathfrak S_4\times\{\pm1\}).
\end{aligned}\right.
\tag*{$(49)$}
\]

Eh bien on suppose \deleted{\unsure{pas}} que $T_0\,3,T_0\,4$, (pas que tous les sommets de = ordre), alors $P_0$ est la double pyramide, $Z'=\{1,\underline-\}$ si $n$ impair ou $n=4$, $Z'\simeq\{1,\sigma\}\times\{1,\underline-\}$ dans le cas $n$ pair $\geqslant5$ $]\!]$
% archive sheet number: 123
\scanpage{125}\mspage{63}

Ainsi, la classe $\Pi$ \inserted{de structures polyédrales ``équivalentes'' sur $S$} se trouve être un espace homogène sous \unsure{le} groupe \uline{commutatif} $Z'$, il en résulte que le stabilisateur d'un $i\in\Pi$ dans $Z'$ est indépendant de $i$, c'est donc aussi le stabilisateur de $F_0$, \inserted{ou de $A_0$} i.e. formé des $\lambda\in Z'=\operatorname{Comm}_{\unsure{\mathfrak S_S|G_S}}(G_0^{\unsure{+|\ast}})$ tels qu'on ait $\forall s,t\in S$

\[
\{s,t\}\in A_0\Longleftrightarrow\{s,\lambda t\}\in A_0
\tag*{$(50)$}
\]

ce qui implique d'ailleurs (remplaçant $\{s,t\}$ par $\{\lambda t,s\}$

\[
\deleted{\unsure{(51)}}\quad
\underbrace{\{\lambda t,\deleted{\unsure{\lambda t|\lambda s}}s\}}_{=\{s,\lambda t\}}\in A_0
\Longleftrightarrow\{\lambda s,\lambda t\}\in A_0
\]

d'où

\[
\lambda\in G_0\quad\text{et même}\quad \lambda\in\operatorname{Cent}(G_0)
\]

puisque $\lambda$ centralise $G_0$, \deleted{Donc} Mais \inserted{(mieux)} on peut dire que (50) signifie aussi

\[
\operatorname{Et}_{A_0}(t)=\operatorname{Et}_{A_0}(\lambda t)
\tag*{$(51)$}
\]

-- ce qui implique que $\lambda=\mathrm{id}$, sauf dans le seul cas où \deleted{\unsure{des|les}} l'étoile ne détermine pas son centre, i.e. le cas \deleted{\unsure{fini|bien}} de l'octaèdre -- \unsure{cas|ou}, si on ne retient pour l'instant que $T_0\,3,T_0\,4$ celui de la double pyramide. On trouve donc

\[
\begin{aligned}
&\deleted{Z'\ \text{groupe de perm. d'autom. de}}\\
&\deleted{\text{Sous les conditions }T_0\,3,T_0\,4,\ \text{si }\unsure{\lambda\in Z'}}\\
&\deleted{\unsure{Z'}\ \text{\unsure{trans.}}}\quad \Pi\text{ est un torseur sous }Z',\ \text{i.e. le stabilisateur}\\
&\text{de }\deleted{\unsure{F_0}}A_0\text{ dans }Z'\text{ est }\{1\},\ \text{i.e. (50) implique}\\
&\lambda=\unsure{1|\mathrm{id}},\ \text{sauf dans le cas }\deleted{\text{de la}}\text{ où }P_0\text{ est une}\\
&\text{double pyramide }(49\,b),\ \text{où on a }Z'\subset\deleted{\unsure{Z}}G_0,\ \deleted{\text{et}}\\
&\text{où le stabilisateur est réduit à }\{1,\underline-\}.\quad\deleted{\unsure{\lambda\text{ si}}}
\end{aligned}\tag*{$(5\unsure{1|2})$}
\]

% slanted
\margin{moyennant $T_0\,3,T_0\,4$}
% archive sheet number: 124
\scanpage{126}

Dans le cas exceptionnel de la double pyramide, on a $\Pi$ réduit à un élément \unsure{$A_0$|$F_0$|$P_0$} \deleted{\unsure{et}} i.e. $Z'=\{1,-\}$ sss $n$ impair ou $n=4$, et dans le cas inverse ($n$ pair $\geqslant6$) $\Pi$ a exactement deux éléments, qui correspondent à la double pyramide sur le = ensemble $E$ ; muni des deux structures polygonales différentes (d'ordre \deleted{pair} \inserted{$\deleted{\unsure{n}}n\equiv0\ (4)$ et} $\geqslant8$) déduites l'une de l'autre par ``multiplication par $\frac{n}{2}-1$''\ldots

\noindent\begin{tabular}{|p{.92\linewidth}}
\unsure{$\vert$|!}\quad Je laisse tomber par la suite ce cas de la double pyramide, sauf dans le cas de l'octaèdre, i.e pratiquement je demande que les \deleted{points}\inserted{sommets} aient tous même ordre -- donc le seul cas exceptionnel où $\Pi$ n'est pas un torseur sous \deleted{\unsure{$Z$}} $Z'$ est alors celui où $P_0$ est l'octaèdre, $\Pi$ réduit à un élément, $Z'=\{1,\underline-\}$. 
\end{tabular}

Notons que par transport de structure, toute permutation $f$ de $S$ qui transforme \uline{un} élément $F_0$ de $\Pi$ en un autre $F_1\in\Pi$, transforme \inserted{\unsure{=}} $\Pi$ \inserted{tout entier} en lui-\deleted{\unsure{=}}même, et induit une bijection de $\Pi$ avec lui =. Soit donc
% archive sheet number: 125
\scanpage{127}\mspage{64}

\[
 G=\left\{f\in\unsure{\mathfrak S_S|G_S}\ \middle|\
 \begin{aligned}
 &f\text{ prolongé à }\unsure{\mathbb P|\mathcal P|\mathfrak P}(\mathbb P_2(S))\\
 &\text{applique }\deleted{\unsure{g}}\ \Pi\subset\mathbb P(\mathbb P_2(S))\\
 &\text{en lui}=\text{, donc satisfait}\\
 &f(\Pi)=\Pi
 \end{aligned}\right\} \tag*{(53)}
\]

% slanted
\margin{\uline{NB} Si $\Pi$ \unsure{s'ident.|réduit|s'ident} à un élément, alors $G=G_0$ et tout ce qui suit est tautologique, \unsure{NB|OBS} donc $\#\text{ de }\Pi\geqslant2$, ce qui implique qu'on n'est pas dans le cas de l'octaèdre, donc que $\Pi$ est bien un \uline{torseur} sous $Z'$.}

\deleted{L'axiome} \deleted{Une} Notons qu'on a manifestement

\[
\begin{aligned}
G_0&\overset{\text{def}}{=}\operatorname{Aut}(S,A_{\deleted{\unsure{i}}i_0})\subset G\\[-2pt]
\shortparallel&\\[-2pt]
&\text{Stabilisateur de }i_0\in\Pi\text{ dans }G,\text{ opérant sur }\Pi\\
&=\text{noyau de }\ G\ \unsure{\text{---}|\longrightarrow}\ \mathfrak S_\Pi
\end{aligned}\tag*{(54)}
\]

donc $G_0$ est un sous-groupe distingué de $G$, et on a

\[
G/G_0\lhook\joinrel\longrightarrow\mathfrak S_\Pi. \tag*{(55)}
\]

Mais $G$ normalisant $G_0$, normalise aussi son centralisateur $Z'$, donc opère sur $Z'$, et $G_0$ y opère trivialement, donc a priori

\[
G/G_0\longrightarrow\operatorname{Aut}_{\mathrm{gr}}(Z'), \tag*{(56)}
\]

et il est clair par transport de structure que (55) est compatible avec la structure de $Z'$-espace homogène \deleted{\unsure{sur}}\inserted{sur $\Pi$} -- en fait de $Z'$-torseur $\Pi$, cf annotation plus haut -- i.e. on a a priori

\[
G/G_0\longrightarrow\operatorname{Aff}_{Z'}(\Pi). \tag*{(57)}
\]
% archive sheet number: 126
\scanpage{128}

Ce serait le moment d'ailleurs de noter que l'on a

\[
Z'\subset G \tag*{(58)}
\]

car on voit \deleted{que} pour $(s,t)\in S^2$

\[
\{s,t\}\in A_{\deleted{\unsure{i}}i_0}\Longleftrightarrow
\{s,\lambda t\}\in A_{(\lambda i_0)}\Longleftrightarrow
\{\lambda s,\lambda t\}\in A_{\underbrace{\lambda.(\lambda i_0)}_{\lambda^2 i_0}}
\]

d'où.

\[
\underbrace{\lambda(A_{i_0})}_{\substack{\text{transport}\\\text{de structure}}}
=A_{\underbrace{\lambda^2.(i_0)}_{\substack{\text{opération}\\\text{de }Z'\text{ sur}\\\Pi\text{ définie par (45)}\\\text{p. }61'}}}, \tag*{(59)}
\]

% slanted
\margin{attention, tournant dangereux et \deleted{\unsure{important}}\inserted{\unsure{mystérieux}} car $\operatorname{det}.(\lambda\,\mathrm{id}_M)=\lambda^2$ (\unsure{$\#$|$\ne$} $\lambda$ en \deleted{\unsure{tant que|général}} que fonction en $\lambda$!).}

et l'opération de $G$ sur $Z'$ est induite par autom. intérieurs. Ceci dit, on va poser

\[
\begin{aligned}
(T_0\,7)\quad&Z'\ \deleted{\text{\unsure{centralise}}}\inserted{\text{est central}}\text{ dans }G\quad\text{(NB. il contient}\\
&\text{de }\deleted{\text{\unsure{toutes}}}\text{ façons le centre de }G,\text{ donc }T_0\,7\\
&\text{signifie aussi}\\
&\hspace{2em}Z'=\operatorname{Cent}(G)\unsure{,|)}
\end{aligned}\tag*{(60)}
\]

qui équivaut aussi au fait que $G$ opère sur $\Pi$ par translations, i.e. (56) trivial i.e. (57) donné par un hom de groupes

\[
\deleted{\unsure{Z}}G\xrightarrow{\delta}Z', \tag*{(61)}
\]

de sorte qu'on a finalement une suite exacte. canonique
% archive sheet number: 127
